\documentclass[11pt]{article}
\usepackage[T1]{fontenc}
\usepackage[utf8]{inputenc}
\usepackage{lmodern}
\usepackage[a4paper,margin=1.1in]{geometry}
\usepackage{microtype}
\PassOptionsToPackage{hyphens}{url}
\usepackage[hidelinks]{hyperref}
\usepackage{parskip}
\setlength{\parskip}{0.7em}

\title{\textbf{Learn robot autonomy at Stanford, now open to everyone}\\[0.3em]
\large Principles of Robot Autonomy (PoRA): textbook, lectures, exercises and teaching materials}
\author{Marco Pavone (\href{https://x.com/drmapavone}{@drmapavone})}
\date{Posted on X, 22 October 2026\\[0.3em]
\small\url{https://x.com/drmapavone/status/2113282303564956047}}

\begin{document}
\maketitle

Ever wondered what it's like to learn robot autonomy @Stanford? Now you can.

Today, we're opening up the robot autonomy teaching materials that my collaborators and I have been developing over the past several years---for anyone, anywhere, completely free.

At the heart of this effort is \emph{Principles of Robot Autonomy} (PoRA), a textbook inspired by the Stanford robot autonomy course series I've been teaching since 2017.

Robot autonomy has evolved from a research discipline into a technology deployed in the real world. Building capable autonomous systems requires more than understanding individual algorithms: it requires understanding how perception, state estimation, planning, control, and learning work together in a complete autonomy stack.

That's what we set out to teach with PoRA: bringing classical robotics and modern AI together in one coherent framework, grounded in both theory and practice.

Coming soon in print from Cambridge University Press, \emph{Principles of Robot Autonomy} is also freely available online and comes with a comprehensive collection of open educational resources:

\begin{description}
  \item[Textbook,] \url{https://arxiv.org/abs/2608.03496} --- Foundations of robot autonomy, from classical methods to modern AI-based approaches.
  \item[Video lectures,] \url{https://porabook.com/lectures/} --- Lectures from my Stanford courses on robot autonomy and on optimal and learning-based control.
  \item[Interactive exercises,] \url{https://porabook.com/exercises/} --- Hands-on Python implementations and Jupyter notebooks.
  \item[Teaching materials,] \url{https://porabook.com/instructors/} --- Lecture slides and resources for instructors.
  \item[Explore all PoRA resources,] \url{https://porabook.com} --- Including Ask the Book, an AI study assistant integrated throughout the textbook to answer questions, explain derivations, and connect ideas across chapters.
  \item[Read the Stanford announcement:] \href{https://news.stanford.edu/stories/2026/10/autonomous-robotics-education-curriculum-textbook}{\nolinkurl{news.stanford.edu/stories/2026/10/}\\ \nolinkurl{autonomous-robotics-education-curriculum-textbook}}
\end{description}

Our goal is simple: to make high-quality education in robot autonomy accessible to students, researchers, engineers, educators, and anyone curious about how autonomous robotic systems work.

I believe that sharing knowledge is fundamental to the mission of a university. Autonomous robotics---and Physical AI more broadly---are among the defining technological frontiers of our time. We want to help more people understand these technologies, build with them, and advance the field.

A huge thank you to my co-authors, @DanieleGammelli, Joseph Lorenzetti, @katielulula, and @GioeleZardini, and to everyone who helped bring this project to life.

Explore the resources, use them in your own learning or teaching, and please share them with anyone who might find them useful!

\textit{\#Robotics \#PhysicalAI \#ArtificialIntelligence \#AutonomousSystems \#Education}

\textit{@Stanford @StanfordEng @StanfordAILab @StanfordHAI}

\vfill
\begin{center}\footnotesize
Text reproduced from the original post on X for easier reading; emoji bullets are set as labels. All words are Marco Pavone's.
\end{center}
\end{document}
